\scanpage{332}
\mspage{31}

A l'opposé des infirmités congénitales de l'égo, un tel état de maladie de la psyché et de l'âme n'est nullement une désolante fatalité. Ce n'est pas notre état lors de notre naissance sur terre, et il ne s'instaure jamais sans qu'il n'y ait de la part de l'âme un acte ou des actes à cet effet (dont la portée le plus souvent reste inconsciente). Ce sont les ``\textit{actes d'infidélité}'', lesquels (en tenant compte de la subreptice immixtion de l'égo partout dans la psyché...) ne sont ni plus, ni moins que les \textit{actes d'abdication} plus ou moins complète de l'âme devant la poussée envahissante de l'égo. Et aux yeux de Dieu, l'acte le plus grave de tous que l'âme puisse commettre, c'est l'acte d'abdication totale. (Et un tel acte est toujours une abdication \textit{au profit de l'égo}.) C'est cela et rien d'autre qu'il faut entendre par ``\textit{suicide spirituel}''.

Parmi les destinataires de la Lettre, il en est plus de vingt qui sont à l'état de cadavres spirituels. Mais en eux, comme en tous les autres, Dieu agira. Il les arrachera malgré eux à leur mort volontaire. Et les morts vivront, et marcheront...

9. L'infidélité est un état de vanité

Ce qui constamment pèse sur l'âme avec plus ou moins d'insistance ou de véhémence et qui constamment la pousse aux actes d'infidélité, ce sont les poussées issues de la \textit{vanité}, ou de la forme ``dure'' de celle-ci : \textit{l'orgueil}, et le \textit{désir de domination} qui très souvent en découle. \textit{La ``maladie d'enfance'' dont il vient d'être question, maladie de l'égo qui sans cesse tend à envahir toute la psyché, n'est autre que la vanité.} Et cette maladie gagne l'âme, l'âme est ``infidèle'', dès l'instant où elle s'abandonne yeux fermés aux pulsions de vanité provenant de l'égo. Et \textit{tous} ses conditionnements sans exception, tant familiaux que culturels et y compris même (sous une forme parfois plus subtile) les conditionnements ``religieux'', poussent l'âme depuis sa plus tendre enfancé dans la voie d'une telle abdication.

Aussi la vanité est-elle, et de très loin, la maladie la plus répandue parmi les hommes. Et la psyché a été réglée par Dieu de telle façon qu'elle ne voit la vanité pratiquement jamais que chez autrui (chose facile, et qui rarement nous avance à grand chose...). Aussi celui qui a franchi le pas de reconnaître enfin cette maladie tenace en lui-même, et qui dans la joie de son coeur se hasarde à parler de sa grande découverte, se voit aussitôt traité avec des airs de commisération condescendante, coupant court aussitôt à toute velléité de communication...

A la bonne heure, certaines lectures m'avaient mis en tête l'idée (consolatrice et des plus convaincantes...) que dès lors qu'on a clairement reconnu la présence envahissante de la vanité dans sa psyché, tout était fait et que ladite vanité allait, sous ce regard scrutateur, s'évanouir comme par enchantement. Cela fait plus de treize ans à présent que je scrute, sans en être débarassé pour autant. Il est vrai qu'aussi longtemps qu'on tient la garce à l'oeil, elle se tient plus ou moins à carreau. (Mais gare aux moments d'inattention !) Ce qui est sûr, c'est qu'en moi elle a cessé de faire la loi. Et c'est beaucoup...

Mais pour tous ceux à qui s'adresse la Lettre, un tout premier pas à franchir pour reconnaître son état d'infidélité et pour y mettre fin, c'est de prendre connaissance de l'action de la vanité dans son être. Des \textit{résistances} occultes d'une force immense, sur le qui vive à toute heure tant dans l'égo que dans l'âme, s'opposent efficacement à une telle prise de

% archive sheet number: 332
